\SourceLeaf{402}%
\noindent Substituting this, we obtain
\[
\begin{aligned}
&2(Pt't'+2Qt'\theta'+R\theta'\theta')\BareRoot{(PR-QQ)}.(df-d\chi)\\
&\quad=t'^3(-QP'+2PQ'-PP'')du\\
&\qquad+t't'\theta'(PR'-3QP''-\frac{2QQ}{P}P'+4QQ')du\\
&\qquad+t'\theta'\theta'(2QR'-\frac{2QQ}{P}P''-\frac{RQ}{P}P'+2RQ'-RP'')du\\
&\qquad+\theta'^3\left(RR'-\frac{QR}{P}P''\right)du.
\end{aligned}
\]
[But the right-hand side is identically]
\[
\begin{aligned}
=(Pt't'+2Qt'\theta'+R\theta'\theta')\{&\left(-\frac QP P'+2Q'-P''\right)t'du\\
&+\left(R'-\frac QP P''\right)\theta'du\}.
\end{aligned}
\]

We therefore have:
\begin{equation}\tag*{[22]}
\begin{aligned}
&2P\BareRoot{(PR-QQ)}(df-d\chi)\\
&\quad=\{(-QP'+2PQ'-PP'')t'+(PR'-QP'')\theta'\}du\\
&\quad=-QdP+(2PQ't'+PR'\theta')du-PP''t'du,
\end{aligned}
\end{equation}
[for which we can also write:]
\begin{equation}\tag*{[23]}
\begin{aligned}
&2\BareRoot{(PR-QQ)}(df-d\chi)\\
&\quad=\frac{-QP'+2PQ'-PP''}{P}dt-\frac{-RR'+2RQ''-QR''}{R}d\theta\\
&\qquad-Q\left(\frac{P''}{P}-2\frac{Q'}{Q}+\frac{R''}{R}\right)d\theta.
\end{aligned}
\end{equation}
\Needspace{8\baselineskip}
\begin{center}\bfseries Remarks.\end{center}
\noindent\hspace*{1.2em}In \S\ 13 of the \textit{Disquisitiones generales circa superficies curvas}, \textsc{Gauss} holds out the prospect of publishing further investigations on the »absolute« properties of curved surfaces, which, like the measure of curvature and geodesic lines, are invariant under bending of the surfaces. The preceding note is to be understood as preparatory work for this. It most likely dates to the period between 1822 and 1825. It is written on separate slips of paper, evidently in great haste. This probably explains why \textsc{Gauss} sometimes calls the »quantities determining the surface« $t,u$ and sometimes $t,\theta$, whereas $t,\theta$ has been used consistently here. It probably also explains why he later uses the letters $\sigma$ and $\psi$, which occur in equations 3 and 9, for the quantities denoted here by $\tau$ (equation 6) and $\chi$ (equation 17). There are also
